% This file is generated from the corresponding Obsidian note.
% Source: Teoricos/GS - Teo3.md
\chapter{Topología cociente}
\label{ch:teo-3}
\section{Topología cociente}
\begin{bookdefinition}{Topología cociente}
Sean $X$ un espacio topológico, $Y$ un conjunto arbitrario y

\[
\pi:X\longrightarrow Y
\]
una función sobreyectiva. La \textbf{topología cociente} sobre $Y$ determinada por $\pi$ se define declarando que

\[
U\subseteq Y \text{ es abierto}
\iff
\pi^{-1}(U)\subseteq X \text{ es abierto}.
\]
\end{bookdefinition}
\begin{bookdefinition}{Función abierta y función cerrada}
Sean $X$ e $Y$ espacios topológicos y $F:X\to Y$ una función. Se dice que $F$ es:

\begin{itemize}
\item \textbf{abierta} si lleva abiertos en abiertos,
\item \textbf{cerrada} si lleva cerrados en cerrados.
\end{itemize}
\end{bookdefinition}
\begin{bookexercise}{Criterio para Hausdorff}
Sea $\pi:X\to Y$ una función sobreyectiva y abierta. Sea $Y$ con la topología cociente. Entonces $Y$ es Hausdorff si y solo si

\[
R:=\{(p,q)\in X\times X:\pi(p)=\pi(q)\}
\]
es cerrado en $X\times X$.

\end{bookexercise}
\begin{bookproof}{}
Hacerla

\bookqed
\end{bookproof}
\phantomsection\label{blk:d48425}
\begin{bookremark}{Funcion no abierta}
Sea $\pi:X\to Y$ una función sobreyectiva y sea $Y$ con la topología cociente. Dar un ejemplo donde $\pi$ no sea abierta.

\end{bookremark}
\begin{bookproof}{}
Tomemos $X=[0,1]$ con la topología usual y consideremos la relación de equivalencia que identifica solo los extremos:

\[
x\sim y \iff x=y \text{ o bien } \{x,y\}=\{0,1\}.
\]
Sea

\[
Y:=X/{\sim}
\]
con la topología cociente, y sea

\[
\pi:X\to Y,\qquad x\mapsto [x]
\]
la proyección natural. Por definición, $\pi$ es sobreyectiva.

Veamos que $\pi$ no es abierta. Consideremos el abierto

\[
U=[0,1/2)\subseteq [0,1].
\]
Entonces $\pi(U)\subseteq Y$ contiene la clase $[0]=[1]$. Además,

\[
\pi^{-1}(\pi(U))=U\cup\{1\}=[0,1/2)\cup\{1\}.
\]
En efecto, todos los puntos de $U$ vuelven a caer en $\pi(U)$, y también $1$ porque $\pi(1)=\pi(0)$ con $0\in U$.

Pero $[0,1/2)\cup\{1\}$ no es abierto en $[0,1]$, ya que $1$ no es punto interior. Luego, por la definición de topología cociente, $\pi(U)$ no es abierto en $Y$.

Por lo tanto, $\pi$ es una función sobreyectiva que define la topología cociente en $Y$, pero no es abierta.

\bookqed
\end{bookproof}
\begin{bookexercise}{Propiedad universal del cociente}
Sean $X$ y $Z$ espacios topológicos, $\pi:X\to Y$ sobreyectiva, y consideremos en $Y$ la topología cociente. Si $F:Y\to Z$ es una función, entonces

\[
F \text{ es continua}\iff F\circ \pi:X\to Z \text{ es continua}.
\]
\end{bookexercise}
\section{El espacio proyectivo real}
\begin{bookdefinition}{Presentación de $\mathbb{RP}^n$ como cociente de la esfera}
Veamos que $\mathbb{R}\mathbb{P}^{n}$ es variedad topologica de dimension $n$

\end{bookdefinition}
\begin{bookproof}{}
\begin{itemize}
\item Definimos $\mathbb{R}\mathbb{P}^{n}$
\end{itemize}
\begin{enumerate}
\item En $S^n$ consideramos la relación de equivalencia
\end{enumerate}
\[
p\sim q\iff q=\pm p.
\]
es decir, identificamos puntos antipodales. 2. Definimos

\[
\mathbb{RP}^n:=S^n/{\sim} = \{[p]:p\in S^n\}.
\]
Sea

\[
\pi:S^n\longrightarrow \mathbb{RP}^n, \qquad p\longmapsto [p].
\]
\begin{enumerate}[start=3]
\item Dotamos a $\mathbb{RP}^n$ de la topología cociente inducida por $\pi$.
\end{enumerate}
\begin{itemize}
\item Veamos que la proyección $\pi$ es abierta.
\end{itemize}
\begin{enumerate}
\item Para esto deberiamos ver que dado $U\subseteq S^{n}$ abierto entonces $\pi(U)$ es abierto, pero por definicion de topologia cociente es lo mismo que ver que $\pi ^{-1}(\pi(U))$ es abierto
\end{enumerate}
\[
\pi^{-1}(\pi(U))=U\cup (-U),
\]
y el lado derecho es abierto en $S^n$. 2. Por lo tanto $\pi(U)$ es abierto y entonces $\pi$ es abierta

\begin{itemize}
\item Veamos Segundo axioma de numerabilidad
\end{itemize}
\begin{enumerate}
\item Como $\pi$ es abierta y sobreyectiva, y $S^n$ satisface $N_2$, se obtiene que $\mathbb{RP}^n$ también satisface $N_2$.
\end{enumerate}
\begin{itemize}
\item Hausdorff
\end{itemize}
\begin{enumerate}
\item Como $\pi$ es abierta, para probar que $\mathbb{RP}^n$ es Hausdorff, usando el \hyperref[blk:d48425]{\textasciicircum{}d48425} basta ver que
\end{enumerate}
\[
R:=\{(p,q)\in S^n\times S^n:\pi(p)=\pi(q)\}
\]
es cerrado. 2. Pero

\[
R=\{(p,q)\in S^n\times S^n:p=q\}\cup \{(p,q)\in S^n\times S^n:p=-q\},
\]
y ambos conjuntos son cerrados en $S^n\times S^n$.

\begin{itemize}
\item Localmente euclideo
\end{itemize}
\begin{enumerate}
\item Tomamos $[p] \in \mathbb{R}\mathbb{P}^{n}$ obviamente existe un $U_i^\pm$ un casquete de $S^n$ tal que $p \in \pi(U_{i}^{\pm})$
\item Ademas notamos que la restricción
\end{enumerate}
\[
\pi|_{U_i^\pm}:U_i^\pm\longrightarrow \pi(U_i^\pm)
\]
es biyectiva. 3. Como $\pi$ es abierta, luego $\pi(U_i^\pm)$ es abierto en $\mathbb{RP}^n$. 4. Ademas como $\pi$ abierta y biyectiva entonces es homeomorfismo 5. Notar que como restringimos el dominio, la inversa tengra restringido el codominio, osea que al volver no me da ambos representantes de la case (como lo haria $\pi$ sin restringir) si no que me da un representante solo. Digamos $\pi|_{U_{i}^{+}}$ es una especie de nueva funcion que al hacer inveras no devuelve dos elementos por cada clase de $\mathbb{R}\mathbb{P}^{n}$ si no que devuelve solo uno, el que esta en $U_{i}^{+}$\\
6. Por lo tanto

\[
\varphi_i^\pm\circ (\pi|_{U_i^\pm})^{-1}:\pi(U_i^\pm)\longrightarrow B^n
\]
es un homeomorfismo.

En consecuencia, $\mathbb{RP}^n$ es una variedad topológica de dimensión $n$.

\bookqed
\end{bookproof}
\begin{bookdefinition}{Otra presentación de $\mathbb{RP}^n$}
Veamos que $\mathbb{R}\mathbb{P}^{n}$ es variedad topologica de otra manera

\end{bookdefinition}
\begin{bookproof}{}
\begin{itemize}
\item Definimos la rel de equivalencia
\end{itemize}
\begin{enumerate}
\item En $\mathbb{R}^{n+1}\setminus\{0\}$ consideramos la relación de equivalencia
\end{enumerate}
\[
p\sim q \iff \exists \lambda\in \mathbb{R}\setminus\{0\}\text{ tal que }q=\lambda p.
\]
\begin{enumerate}[start=2]
\item Es decir, dos vectores son equivalentes si generan la misma recta que pasa por el origen.
\end{enumerate}
\begin{itemize}
\item Definimos el cociente
\end{itemize}
\begin{enumerate}
\item Definimos
\end{enumerate}
\[
\mathbb{RP}^n:= (\mathbb{R}^{n+1}\setminus\{0\})/{\sim}= \{[p]:p\in \mathbb{R}^{n+1}\setminus\{0\}\}
\]
con la topología cociente determinada por la proyección natural

\[
\pi:\mathbb{R}^{n+1}\setminus\{0\}\longrightarrow \mathbb{RP}^n, \qquad p\longmapsto [p].
\]
\begin{itemize}
\item Hausdorff
\end{itemize}
\begin{enumerate}
\item Si $U\subseteq \mathbb{R}^{n+1}\setminus\{0\}$ es abierto, entonces
\end{enumerate}
\[
\pi^{-1}(\pi(U))= \bigcup_{\lambda\in \mathbb{R}\setminus\{0\}} \lambda U.
\]
\begin{enumerate}[start=2]
\item Como la multiplicación por $\lambda\neq 0$ es un homeomorfismo de $\mathbb{R}^{n+1}\setminus\{0\}$, cada $\lambda U$ es abierto
\item Por lo tanto, $\pi^{-1}(\pi(U))$ es abierto y $\pi(U)$ es abierto. Osea $\pi$ es abierta
\item Luego consideremos
\end{enumerate}
\[
R:=\{(p,q)\in (\mathbb{R}^{n+1}\setminus\{0\})^2:q=\lambda p \text{ para algún } \lambda\neq 0\}.
\]
osea el conjunto de los elementos ld 5. Luego recordamos que dos elementos son ld si todas de su submatrices tienen determinante 0 6. Para $1\le i<j\le n+1$, definimos

\[
Q_{ij}:\bigl(\mathbb{R}^{n+1}\setminus\{0\}\bigr)^2\longrightarrow \mathbb{R},\qquad Q_{ij}(x,y):= \det\begin{pmatrix}x_i & y_i\\ x_j & y_j\end{pmatrix}.
\]
\begin{enumerate}[start=7]
\item Entonces
\end{enumerate}
\[
R=\bigcap_{1\le i<j\le n+1} Q_{ij}^{-1}(\{0\}),
\]
luego $R$ es cerrado. Por ser interseccion finita de cerrados ($Q_{ij}$ escontinua y $\{ 0 \}$ cerrado) 8. Por consiguiente, $\mathbb{RP}^n$ es Hausdorff.

\begin{itemize}
\item Localmente euclideo
\end{itemize}
\begin{enumerate}
\item Sea $[p]\in \mathbb{RP}^n$ y tomemos un representante
\end{enumerate}
\[
p=(x_1,\dots,x_{n+1})
\]
con alguna coordenada $x_i\neq 0$. 2. Definimos

\[
\widetilde U_i:=\{(y_1,\dots,y_{n+1})\in \mathbb{R}^{n+1}\setminus\{0\}:y_i\neq 0\}
\]
que es abierto en $\mathbb{R}^{n+1}\setminus\{0\}$ 3. Luego

\[
U_i:=\pi(\widetilde U_i)
\]
que es un abierto de $\mathbb{RP}^n$ por que en este caso $\pi ^{-1}(U_{i})=\tilde{U}_{i}$ que es abierto (osea justo $\tilde{U}_{i}$ es saturado) 4. La aplicación

\[
\varphi_i:U_i\longrightarrow \mathbb{R}^n, \qquad [(x_1,\dots,x_{n+1})] \longmapsto \left( \frac{x_1}{x_i},\dots,\widehat{\frac{x_i}{x_i}},\dots,\frac{x_{n+1}}{x_i} \right)
\]
está bien definida y es un homeomorfismo. 5. Una inversa explícita está dada por

\[
\psi_i(u_1,\dots,u_n) = [u_1,\dots,u_{i-1},1,u_i,\dots,u_n].
\]
\begin{enumerate}[start=6]
\item Por lo tanto, los abiertos $U_i$ dan una familia de cartas de $\mathbb{RP}^n$ con imagen en $\mathbb{R}^n$.
\end{enumerate}
\bookqed
\end{bookproof}
\section{El toro}
\begin{bookexercise}{Una presentación del toro REVISAR}
Ver a $\mathbb{R}^n$ como grupo abeliano con la suma. Entonces $\mathbb{Z}^n$ es subgrupo de $\mathbb{R}^n$. Se denota por

\[
\mathbb{T}^n:=\mathbb{R}^n/\mathbb{Z}^n
\]
el conjunto de clases laterales a izquierda de $\mathbb{Z}^n$. Equivalentemente,

\[
x\sim y \iff x-y\in \mathbb{Z}^n.
\]
Dotamos a $\mathbb{T}^n$ de la topología cociente inducida por la proyección natural

\[
\pi:\mathbb{R}^n\longrightarrow \mathbb{T}^n,\qquad x\longmapsto [x].
\]
Ver que $\mathbb{T}^n$ es una variedad topológica de dimensión $n$.

\end{bookexercise}
\begin{bookproof}{}
\begin{itemize}
\item Definimos la relación de equivalencia
\end{itemize}
\begin{enumerate}
\item Para $x,y\in \mathbb{R}^n$, tenemos
\end{enumerate}
\[
x\sim y \iff x-y\in \mathbb{Z}^n.
\]
\begin{enumerate}[start=2]
\item Es decir, $x$ e $y$ son equivalentes si difieren en un vector con coordenadas enteras.
\end{enumerate}
\begin{itemize}
\item Veamos que la proyección $\pi$ es abierta
\end{itemize}
\begin{enumerate}
\item Sea $U\subseteq \mathbb{R}^n$ abierto. Entonces
\end{enumerate}
\[
\pi^{-1}(\pi(U))=\bigcup_{m\in \mathbb{Z}^n}(U+m).
\]
\begin{enumerate}[start=2]
\item Como para cada $m\in \mathbb{Z}^n$ la traslación $x\mapsto x+m$ es un homeomorfismo de $\mathbb{R}^n$, cada $U+m$ es abierto.
\item Luego $\pi^{-1}(\pi(U))$ es abierto, y por definición de topología cociente $\pi(U)$ es abierto.
\item Por lo tanto, $\pi$ es abierta.
\end{enumerate}
\begin{itemize}
\item Segundo axioma de numerabilidad
\end{itemize}
\begin{enumerate}
\item Como $\pi$ es abierta y sobreyectiva, y $\mathbb{R}^n$ satisface $N_2$, se obtiene que $\mathbb{T}^n$ también satisface $N_2$.
\end{enumerate}
\begin{itemize}
\item Hausdorff
\end{itemize}
\begin{enumerate}
\item Como $\pi$ es abierta, usando el \hyperref[blk:d48425]{\textasciicircum{}d48425} basta ver que
\end{enumerate}
\[
R:=\{(x,y)\in \mathbb{R}^n\times \mathbb{R}^n:\pi(x)=\pi(y)\}
\]
es cerrado. 2. Pero

\[
R=\{(x,y)\in \mathbb{R}^n\times \mathbb{R}^n:x-y\in \mathbb{Z}^n\}
=\bigcup_{m\in \mathbb{Z}^n}\{(x,y)\in \mathbb{R}^n\times \mathbb{R}^n:y=x-m\}.
\]
\begin{enumerate}[start=3]
\item Cada conjunto
\end{enumerate}
\[
G_m:=\{(x,y)\in \mathbb{R}^n\times \mathbb{R}^n:y=x-m\}
\]
es cerrado, pues es el conjunto de ceros de la función continua $(x,y)\mapsto y-x+m$. 4. Además, la unión $\bigcup_{m\in \mathbb{Z}^n}G_m$ es cerrada: si $(x_k,y_k)\in R$ y $(x_k,y_k)\to (x,y)$, entonces $y_k-x_k\in \mathbb{Z}^n$ para todo $k$, y como $y_k-x_k\to y-x$ y $\mathbb{Z}^n$ es cerrado en $\mathbb{R}^n$, se sigue que $y-x\in \mathbb{Z}^n$. 5. Luego $R$ es cerrado y por consiguiente $\mathbb{T}^n$ es Hausdorff.

\begin{itemize}
\item Localmente euclídeo
\end{itemize}
\begin{enumerate}
\item Sea $[x]\in \mathbb{T}^n$. Tomamos $\varepsilon<1/2$ y sea $B(x,\varepsilon)\subseteq \mathbb{R}^n$ la bola abierta de centro $x$ y radio $\varepsilon$.
\item Si $m\in \mathbb{Z}^n\setminus\{0\}$, entonces $\|m\|\ge 1$. Por lo tanto, $B(x,\varepsilon)\cap (B(x,\varepsilon)+m)=\varnothing$.
\item En consecuencia, la restricción
\end{enumerate}
\[
\pi|_{B(x,\varepsilon)}:B(x,\varepsilon)\longrightarrow \pi(B(x,\varepsilon))
\]
es biyectiva. 4. Como $\pi$ es abierta, $\pi(B(x,\varepsilon))$ es abierto en $\mathbb{T}^n$. Entonces $\pi|_{B(x,\varepsilon)}$ es un homeomorfismo. 5. Como $B(x,\varepsilon)$ es homeomorfo a una bola abierta de $\mathbb{R}^n$, resulta que cada punto de $\mathbb{T}^n$ tiene un entorno abierto homeomorfo a un abierto de $\mathbb{R}^n$.

En consecuencia, $\mathbb{T}^n$ es una variedad topológica de dimensión $n$.

\bookqed
\end{bookproof}
